% F5e -- What top-1, top-3 and top-5 mean, and how a score is computed from them.
\begin{figure}[!htbp]
\centering
\fitwidth{%
\begin{tikzpicture}[font=\small,
  cand/.style={draw=clrule,fill=white,rounded corners=2pt,align=center,
               inner sep=4pt,text width=2.15cm,font=\scriptsize},
  hit/.style={cand,draw=clgain,fill=clgain!14,thick},
  brk/.style={clmodel,thick,|-|},
  cell/.style={draw=clrule,minimum width=0.85cm,minimum height=0.55cm,
               inner sep=0pt,font=\scriptsize},
  ok/.style={cell,draw=clgain,fill=clgain!40},
  ko/.style={cell,draw=clnonsig,fill=white},
]
% ================= (a) one prediction, ranked =================
\node[anchor=west,font=\footnotesize\bfseries,text=clmodel] at (-0.35,3.55)
  {(a) The model never answers once. It ranks all 369 antennas.};
\node[anchor=west,font=\scriptsize,text=clrule] at (-0.35,3.18)
  {Here are its first five, for one moment of one person's day, with how sure it
   is of each.};

% the three depth markers, stacked above the row
\draw[brk] (-0.05,2.62) -- (2.35,2.62);
\node[font=\scriptsize,text=clmodel,anchor=west] at (2.55,2.62) {\textbf{top-1}};
\draw[brk] (-0.05,2.20) -- (7.45,2.20);
\node[font=\scriptsize,text=clmodel,anchor=west] at (7.65,2.20) {\textbf{top-3}};
\draw[brk] (-0.05,1.78) -- (12.55,1.78);
\node[font=\scriptsize,text=clmodel,anchor=west] at (12.75,1.78) {\textbf{top-5}};

\node[cand] (k1) at (1.15,0.95) {\#1\\ \cid{UKVBER101}\\ 52\,\%};
\node[cand] (k2) at (3.70,0.95) {\#2\\ \cid{UKVTGM104}\\ 23\,\%};
\node[hit]  (k3) at (6.25,0.95) {\#3\\ \cid{BKVPAN1}\\ 13\,\%};
\node[cand] (k4) at (8.80,0.95) {\#4\\ \cid{UKVDVO104}\\ 7\,\%};
\node[cand] (k5) at (11.35,0.95){\#5\\ \cid{BKVBOH1}\\ 5\,\%};
\node[font=\scriptsize,text=clrule,anchor=west] at (12.75,0.95)
  {\dots\ 364 more};

\node[font=\scriptsize,text=clgain,anchor=north,align=center] at (6.25,0.28)
  {$\uparrow$ this is the antenna the person really visited};

\node[font=\scriptsize,anchor=west] at (-0.35,-0.55)
  {\textcolor{clloss}{\textbf{top-1: wrong}}, it is not first. \qquad
   \textcolor{clgain}{\textbf{top-3: right}}, it is in the first three. \qquad
   \textcolor{clgain}{\textbf{top-5: right}}, it is in the first five.};

% ================= (b) how a score is computed =================
\draw[clrule,dotted] (-0.45,-1.05) -- (14.4,-1.05);
\node[anchor=west,font=\footnotesize\bfseries,text=clmodel] at (-0.35,-1.45)
  {(b) The score is then just a count, over every graded moment of every test
   person.};

\foreach \i/\st in {0/ok,1/ko,2/ok,3/ok,4/ko,5/ok,6/ok,7/ok,8/ko,9/ok}{
  \node[\st] at (0.1+\i*0.92,-2.25) {};
}
\node[font=\scriptsize,anchor=west,text width=6.2cm] at (9.4,-2.25)
  {shaded = right, empty = wrong\\[2pt]
   7 right out of 10 graded moments\\ $\Rightarrow$ a score of \textbf{70\,\%}};
\node[font=\scriptsize,text=clrule,anchor=west] at (-0.35,-2.95)
  {The one-line rule is graded on the very same moments, and its own count is
   what my score is compared against.};
\end{tikzpicture}}
\caption{Top-1, top-3 and top-5, and the way a score is obtained}
\label{fig:topk}
\figtext{%
\figpara{Panel (a): the model does not give an answer, it gives an order.}{Asked
what comes next, it assigns a degree of confidence to all 369 antennas and sorts
them, most likely first: here 52\,\% for \cid{UKVBER101}, 23\,\% for
\cid{UKVTGM104}, and so on. \textbf{Top-1} asks whether the right answer is in
first place, \textbf{top-3} whether it is anywhere in the first three,
\textbf{top-5} in the first five. In the example the right answer sits third, so
the same prediction fails at top-1 and succeeds at the other two. Reporting three
depths is the convention of this literature rather than a choice of mine
\cite{luca2021survey,qin2025mobllm}.}

\figpara{Panel (b): where a percentage comes from.}{The score is a plain count.
Every graded moment is answered, marked right or wrong at the chosen depth, and
the score is the number of right answers over the number of moments: ten moments,
seven right, 70\,\%. The one-line rule answers the very same moments and gets its
own count, and every result in this report is the difference between the two.}}

\end{figure}
